% =====================================================================
%  Spending Smarter in SACU — Section 3 (LaTeX replication)
%  Source: docs/2026-06_WP-Alexis/Spending-Smarter-in-SACU.docx
%  Self-contained, compilable with pdflatex.
% =====================================================================
\documentclass[11pt]{article}

\usepackage[T1]{fontenc}
\usepackage[a4paper,margin=2.5cm]{geometry}
\usepackage{graphicx}
\usepackage{booktabs}
\usepackage{caption}
\usepackage{amsmath}
\usepackage{siunitx}
\usepackage{xcolor}
% Cross-references (Figure/Table) and citations shown as clickable blue links.
\definecolor{linkblue}{HTML}{1E88E5}
\usepackage[colorlinks=true,linkcolor=linkblue,citecolor=linkblue,urlcolor=linkblue]{hyperref}

% Bibliography: biblatex with Chicago author-date (matches the SDN reference
% format). Citations via \textcite / \parencite are colored by citecolor above.
\usepackage{csquotes}
\usepackage[authordate,backend=biber]{biblatex-chicago}
\addbibresource{refs.bib}

% --- light styling so captions/sources/notes read like the SDN ----------
% Caption + subcaption sit on TOP of each float and are centered.
\captionsetup{font=small,labelfont=bf,justification=centering,singlelinecheck=false}
\graphicspath{{figures/}{../chartReplication/}}

% Subcaption: centered, directly under the caption (top of the float).
\newcommand{\subcap}[1]{{\footnotesize #1\par}}

% Source / Note: left-aligned, at the BOTTOM of the float.
\newcommand{\figsource}[1]{\par\vspace{4pt}{\raggedright\footnotesize\textsc{Source:} #1\par}}
\newcommand{\fignote}[1]{{\raggedright\footnotesize\textsc{Note:} #1\par}}


% Efficiency-gap symbols (kept as macros so they are easy to restyle later)
\newcommand{\eGI}{\ensuremath{e_{GI}}}
\newcommand{\eGH}{\ensuremath{e_{GH}}}
\newcommand{\eGHedu}{\ensuremath{e_{GH}^{\,edu}}}
\newcommand{\eGHhlt}{\ensuremath{e_{GH}^{\,hlt}}}
\newcommand{\eGHblend}{\ensuremath{e_{GH}^{\,blend}}}

\begin{document}

% --- Introduction (excerpt) ------------------------------------------------
% The full Introduction lives in the Word draft; this carries the relocated
% "efficiency gaps as a wedge" paragraph moved out of former 3.1 (revComments R2).
\setcounter{section}{0}
\section{Introduction}
\label{sec:intro}

At the heart of the paper's quantification, the estimated efficiency gaps enter
directly as a wedge on public spending: an efficiency gap of, say, 0.35 implies
that 35 percent of each unit spent does not reach the productive public capital
stock---lost to over-invoicing, poor project selection, weak implementation, or
inadequate maintenance. The appeal of this formulation is that it separates the
level of spending from its productivity---a country spending as much as its
peers but with a larger gap wastes a higher share---so that closing the gap
raises output without additional budget resources.

% --- Section 3 -------------------------------------------------------------
% Mirror the working paper's numbering: Section 3, Figures 3--4, Table 1.
\setcounter{section}{2}
\setcounter{figure}{2}

\section{Measuring Spending Efficiency: Methodology and SACU Results}
\label{sec:efficiency}

This section explains the stochastic-frontier methodology used in the Fiscal
Monitor \parencite{imf2025fm}, reports the SACU-specific efficiency-gap
estimates that enter the paper's simulations, and discusses interpretation and
caveats.

% ---------------------------------------------------------------------
\subsection{Methodology}
\label{sec:methodology}

This paper relies on the spending-efficiency measures established in
\textcite{bankowskichosher2026}, which underpin the October 2025 Fiscal
Monitor \parencite{imf2025fm}. Their approach compares public spending (the input) with
the socioeconomic outcomes it is meant to deliver (the output) in each of the
four spending sectors covered by the framework---public investment, education,
health, and research and development.\footnote{R\&D lies outside the scope of
this paper and is not considered further; the analysis focuses on public
investment, education, and health.} Stochastic frontier analysis (SFA) estimates
a best-practice production function from a cross-country panel and measures each
country's efficiency as its distance below that frontier. Unlike non-parametric
frontiers such as data envelopment analysis, SFA separates genuine inefficiency
from statistical noise. It also handles persistent structural differences
between countries---such as being landlocked or having a large private health
sector. These are absorbed separately, through the four-component (``true random
effects'') estimator of \textcite{filippinigreene2016}. The estimator
distinguishes time-invariant structural heterogeneity from persistent
inefficiency, so the former is not mislabelled as the latter.

Inputs are measured as a five-year moving average of real public spending per
capita in purchasing-power-parity terms; outputs combine multiple quantity and
quality indicators for each sector (for public investment, for example, road and
electricity access alongside infrastructure-quality ratings; for education,
enrolment, completion, and learning measures; for health, life expectancy,
immunization, and health-system capacity). To avoid sensitivity to any single
indicator, scores are averaged across all combinations of the available
indicators. Because the estimator is time-varying, it delivers a separate score
for each individual country in each individual year, rather than a single
cross-sectional snapshot. Each score lies between 0 and 1, with 1 denoting a
country on the frontier. Throughout this
paper we work with the complementary \emph{efficiency gap}---one minus the
score---so that 0 denotes full efficiency and higher values denote larger
shortfalls. Full methodological detail is given in
\textcite{bankowskichosher2026} and the Fiscal Monitor Online Annex
\parencite{imf2025annex}.

% ---------------------------------------------------------------------
\subsection{SACU Results}
\label{sec:sacu-results}

Efficiency gaps persist across sectors in SACU. Figure~\ref{fig:eff-gaps}
reports the Fiscal Monitor efficiency gaps for the five SACU countries in three
sectors---infrastructure, health, and education---at two points in time (2000
and 2023), alongside the unweighted averages for advanced economies, emerging
markets, and sub-Saharan Africa. Three patterns stand out. In every sector,
SACU gaps sit above the advanced-economy average and close to the
emerging-market level. Infrastructure has improved the most since 2000. Health
gaps remain the largest and most persistent, while education has stayed broadly
flat even when it deteriorated across advanced economies and emerging markets.

\textit{Infrastructure} gaps have narrowed but remain wider than peers. The
left panel of Figure~\ref{fig:eff-gaps} shows that all five SACU countries
closed their infrastructure efficiency gaps substantially between 2000 and
2023. Even so, the SACU average (0.39 in 2023) remains above the
advanced-economy average (0.36) and close to the emerging-market level (0.37). The progress
reflects both rising spending and strengthened investment-management
frameworks, but the remaining gap points to scope for further gains.

\textit{Health} gaps are persistent and larger than peers. The middle panel shows
health efficiency gaps of 0.36 to 0.49 for SACU countries, little changed since
2000 (when the average gap was 0.42 compared to 0.41 in 2023), and above the peer-group averages.
Health outcomes in SACU continue to reflect the lingering effects of the HIV/AIDS epidemic,
but also structural features of the health systems---urban-rural coverage gaps, workforce
pressures, and supply-chain weaknesses---that limit how efficiently spending
translates into outcomes.

\textit{Education} gaps have been broadly stable. The right panel shows
education efficiency gaps of 0.31 to 0.44 across SACU. Consequently, the
SACU average gap (0.35 in 2023) is in line with the emerging-market average
(0.35) and below the sub-Saharan Africa average (0.37). Most countries
experienced limited change in efficiency between 2000 and 2023. Lesotho
stands out as the exception, with its gap narrowing notably from 0.44 to 0.33.
The gaps in South Africa and Namibia ended modestly above their
2000 levels, reaching 0.41 and 0.37 respectively.

\begin{figure}[htbp]
  \centering
  \caption{Efficiency Gaps in Health, Education, and Infrastructure Spending}
  \label{fig:eff-gaps}
  \subcap{(Comparison between 2000 and 2023)}
  \vspace{4pt}
  \includegraphics[width=\textwidth]{scoreEvolution_SACU_eff_gap_snapshot.png}
  \figsource{IMF staff estimates based on \textcite{imf2025fm}.}
  \fignote{The chart compares 2000 and 2023 values across three sectors. The thick dark-grey line is the SACU aggregate; thin solid lines show individual countries; dashed/dotted lines show reference averages. The y-axis range is 0.2--0.6; Lesotho's off-chart 2000 value (0.88) is annotated.}
\end{figure}

Spending-efficiency differences appear small but place SACU countries far
apart in international comparison. Figure~\ref{fig:infra-dist} sets these
gaps in global context. It plots the 2023 efficiency gaps for every country in
the Fiscal Monitor database across the three sectors, ordered from
highest to lowest gap, with the five SACU countries highlighted. In infrastructure (top panel), Namibia,
Lesotho, and Eswatini sit close to the median of the global distribution;
Botswana and South Africa are modestly below the median. None of the five SACU
countries is among the best performers, confirming that the SACU distribution
represents a set of economies with meaningful but not extreme scope for
improvement.

\begin{figure}[htbp]
  \centering
  \caption{Spending-Efficiency Gaps Across Countries, 2023}
  \label{fig:infra-dist}
  \subcap{Cross-country distribution of spending-efficiency gaps in three sectors,
  with SACU countries highlighted.}
  \vspace{4pt}
  \includegraphics[width=\textwidth]{scoreEvolution_SACU_infra_dist.png}
  \figsource{IMF staff estimates based on \textcite{imf2025fm}.}
  \fignote{Each bar is one country; SACU members highlighted in blue.}
\end{figure}

Table~\ref{tab:eff-gaps} summarizes the country-specific efficiency gaps that
feed the simulation exercise in Sections 5 and 6. The infrastructure gap is
measured for the latest available year (2023). For human capital, we report
both the education-specific gap (latest year 2024) and the health-specific gap
(2023), as well as a blended gap defined as the simple average of the two. The
blended gap is used in the model because the Fiscal Monitor DSGE framework
treats education and health as a single public human-capital stock (schools and
hospitals). The education-specific gap is reported separately because one of the
reallocation variants targets education-sector investment only.

\begin{table}[htbp]
  \centering
  \caption{SACU Efficiency Gaps, Latest Available Year}
  \label{tab:eff-gaps}
  \subcap{Efficiency gap, 0 to 1 scale; higher values indicate larger gaps.
  Used as inputs to the DSGE simulations.}
  \vspace{6pt}
  \begin{tabular}{lcccc}
    \toprule
    Country & Infrastructure & Education & Health & Blended HC \\
            & (\eGI, 2023)   & (\eGHedu, 2024) & (\eGHhlt, 2023) & (\eGHblend) \\
    \midrule
    Botswana     & 0.34 & 0.38 & 0.40 & 0.39 \\
    Eswatini     & 0.39 & 0.35 & 0.41 & 0.38 \\
    Lesotho      & 0.40 & 0.42 & 0.47 & 0.45 \\
    Namibia      & 0.41 & 0.35 & 0.37 & 0.36 \\
    South Africa & 0.33 & 0.37 & 0.37 & 0.37 \\
    \bottomrule
  \end{tabular}
  \figsource{IMF Fiscal Monitor (2025) efficiency-score database, extracted for
  the FM2025\_SACU\_DSGE model.}
  \fignote{Blended human-capital gap is the simple average of the education and
  health gaps.}
\end{table}

% ---------------------------------------------------------------------
\subsection{Interpretation and Caveats}
\label{sec:caveats}

Three interpretive caveats apply:

First, efficiency gaps conflate several underlying factors---institutional
capacity, governance quality, the project mix, and country-specific
shocks---and do not, on their own, identify the mechanism by which resources are
lost. Section 4 links the measured gaps to institutional evidence from the
Public Investment Management Assessments to address this.

Second, estimates are sensitive to the frontier specification. The Fiscal
Monitor Online Annex \parencite{imf2025annex} reports that alternative weightings of the
underlying efficiency-score models can shift the level of estimated gaps by up
to 26 percentage points in the extreme case. We use the baseline specification
throughout, with sensitivity analysis discussed in Section 6.

Third, there are data-window limitations. The SFA uses a five-year moving
average, which means the 2023 estimate reflects spending and outcomes between
2019 and 2023---a period disrupted by the COVID-19 shock, the subsequent SACU
revenue boom, and commodity-price swings that affected Botswana in particular.
Estimates should therefore be interpreted as reflecting the medium-run
institutional capacity of each system rather than a precise snapshot of any
single year.

% --- Appendix --------------------------------------------------------------
\newpage
\appendix
\section{Appendix: Spending-Efficiency Gaps over Time}
\label{sec:appendix}

\begin{figure}[htbp]
  \centering
  \caption{SACU Spending-Efficiency Gaps over Time}
  \label{fig:eff-gaps-evolution}
  \subcap{(Efficiency gap, 0--1 scale; higher values indicate larger gaps)}
  \vspace{4pt}
  \includegraphics[width=\textwidth]{scoreEvolution_SACU_eff_gap.png}
  \figsource{IMF staff estimates based on \textcite{imf2025fm}.}
  \fignote{The thick dark-grey line is the SACU aggregate (unweighted mean of
  the five members); thin solid lines show the individual countries; dotted
  lines show reference averages---advanced economies (AEs), emerging markets
  (EMs), and sub-Saharan Africa (SSA). Hollow circles mark 2000 and 2023.}
\end{figure}

% --- References ------------------------------------------------------------
% IMF (2026) is cited only in figure/table source lines, so \nocite keeps it in
% the list.
\nocite{imf2026sacu}
\printbibliography[title={References}]

\end{document}
